\section{Spatial Support and Emergence of a Macroscopic Newtonian Response}

  \subsection{Discrete Model and Numerical Protocol}

    To illustrate the mechanism described above, we consider a finite
    three-dimensional lattice approximation of a Laplace-type operator.
    Let $L$ denote the discrete graph Laplacian on a cubic grid
    with Dirichlet boundary conditions.
    We define
    \begin{equation}
      A = L + \varepsilon I,
    \end{equation}
    with a small $\varepsilon > 0$ introduced to remove the zero mode
    and ensure invertibility.

    A localized perturbation is implemented by symmetrically modifying
    the weights of edges inside a ball of radius $R$ around a central site.
    The perturbed operator reads
    \begin{equation}
      A \longrightarrow A + \delta A,
    \end{equation}
    where $\delta A$ has compact spatial support.

    The Green column associated with a point source is obtained by solving
    \begin{equation}
      A \phi = \delta_{\mathrm{source}},
    \end{equation}
    and the radial average of $\phi$ is fitted at large distance by
    \begin{equation}
      \phi(r) \sim \frac{a}{r} + b.
    \end{equation}
    The coefficient $a$ measures the amplitude of the effective
    Newtonian response.

    The entropy variation is evaluated locally through the
    support-restricted identity
    \begin{equation}
      \Delta \log \det{}' A
      =
      \log \det (I + (A^{-1})_{SS} \delta A_{SS}),
    \end{equation}
    where $S$ denotes the support of the perturbation.
    In the weak-perturbation regime, this reduces to
    \begin{equation}
      \Delta \log \det{}' A
      \approx
      \mathrm{Tr}((A^{-1})_{SS} \delta A_{SS}),
    \end{equation}
    which was numerically verified.

  \subsection{Support Threshold and Long-Range Response}

    The numerical results reveal a clear distinction between
    purely local entropy variation and macroscopic Newtonian response.

    For perturbations confined to a small spatial region
    ($R=1$ in lattice units), the entropy variation
    $\Delta \log \det{}' A$ is non-zero,
    yet the fitted coefficient $a$ remains essentially unchanged.
    In this regime, the modification of the operator affects
    the spectrum locally without generating a measurable
    long-range shift in the Green amplitude.

    When the spatial support is increased beyond a finite threshold
    ($R \ge 3$ in the simulations), a coherent modification of the
    coefficient $a$ emerges.
    The amplitude shift $\Delta a$ grows monotonically
    with $-\Delta \log \det{}' A$.
    The response becomes macroscopic and consistent with
    a Newtonian $1/r$ behavior.

    These observations indicate that long-range interaction
    does not arise from entropy variation alone,
    but from entropy variation with sufficiently extended
    spatial support.
    The resolvent structure of the operator mediates
    the non-local response, while the extent of the perturbation
    controls whether the effect remains confined or propagates
    to large distances.
